\section[ПРОГРАММНАЯ РЕАЛИЗАЦИЯ И РЕЗУЛЬТАТЫ]{Программная реализация и результаты}
\thispagestyle{plain}

В настоящей главе описаны реализованные функциональные модули
сервиса EduQuiz, представлена инфраструктура его развёртывания,
приведены результаты экспериментальной оценки производительности
методом нагрузочного тестирования и сформулированы методика
педагогического эксперимента и перспективы дальнейшей работы.
Скриншоты пользовательских интерфейсов получены в среде Flutter Web
с использованием демонстрационных учётных записей.

\subsection{Реализованные функциональные модули}

В рамках работы реализованы и работают в режиме «end-to-end»
следующие функциональные модули сервиса:
\begin{enumerate}
    \item модуль аутентификации: регистрация, вход, обновление
          токенов доступа, разграничение доступа по ролям;
    \item модуль преподавателя: управление курсами и банками
          вопросов, редактор вопросов нескольких типов, импорт
          банка вопросов из формата Moodle, психометрическая
          аналитика по банку вопросов;
    \item модуль проведения интерактивной учебной сессии в трёх
          режимах (\texttt{classic}, \texttt{timer}, \texttt{race})
          с настройкой длительности таймера, перемешивания вопросов
          и автоматического перехода между вопросами;
    \item модуль обучающегося: подключение к сессии по числовому
          коду, прохождение сессии в режиме реального времени,
          индивидуальное интервальное повторение по алгоритму
          SuperMemo SM-2, профиль с игровой прогрессией.
\end{enumerate}

В соответствии со структурой клиентского приложения далее приведено
описание интерфейсов и ключевых сценариев каждого модуля. Полный
исходный код реализованного сервиса опубликован в репозиториях
\texttt{eduquiz-backend} (серверная часть, язык Go) и
\texttt{eduquiz-flutter} (кроссплатформенный клиент, Flutter
и Dart); фрагменты реализации ключевых компонентов приведены
в Приложении 1. Количественные характеристики реализованной кодовой
базы приведены в табл.~\ref{tab:codebase-stats}.

\begin{table}[!htb]
\fontsize{11pt}{13pt}\selectfont
\centering
\caption{Количественные характеристики реализованной кодовой базы сервиса EduQuiz}
\label{tab:codebase-stats}
\smallskip
\begin{tabular}{|l|c|c|}
    \hline
    \textbf{Характеристика} & \textbf{Серверная часть} & \textbf{Клиентское приложение} \\
    \hline
        Язык программирования      & Go 1.25     & Dart 3 / Flutter 3.35 \\
    \hline
        Объём кода, строк          & $\approx 12\,000$ & $\approx 18\,000$ \\
    \hline
        Число пакетов / модулей    & 14          & 11 \\
    \hline
        Число эндпоинтов API       & 41          & ---     \\
    \hline
        Число экранов              & ---         & 24      \\
    \hline
        Число миграций схемы БД    & 10          & ---     \\
    \hline
        Число таблиц БД            & 14          & ---     \\
    \hline
\end{tabular}
\end{table}
\FloatBarrier

Серверная часть структурирована по слоям: пакет \texttt{cmd}
содержит точки входа основного приложения и инструмента нагрузочного
тестирования; пакет \texttt{internal/handlers}~--- обработчики
REST-эндпоинтов; \texttt{internal/services}~--- бизнес-логика
(алгоритм SM-2, начисление очков опыта, психометрический
анализатор, парсеры Moodle XML и Moodle GIFT); \texttt{internal/ws}~---
WebSocket-хаб; \texttt{internal/server}~--- инфраструктура
HTTP-сервера, маршрутизация и middleware; \texttt{internal/models}~---
структуры предметной области; \texttt{migrations}~--- миграции
схемы базы данных. Клиентское приложение организовано по принципу
feature-first: каждая ключевая функция вынесена в отдельный модуль
(\texttt{auth}, \texttt{teacher}, \texttt{student}, \texttt{room},
\texttt{profile}, \texttt{review}), содержащий собственные экраны,
провайдеры состояния и сервисы взаимодействия с серверной частью.

\subsection{Модуль аутентификации}

При первом запуске клиентского приложения предусмотрен короткий
вводный сценарий, кратко представляющий основные сценарии
использования сервиса. Признак прохождения вводного сценария
сохраняется в локальном хранилище и повторно сценарий
не отображается.

Экран аутентификации совмещает формы регистрации и входа
(см.~Рисунок~\ref{fig:ui-auth} в Приложении~2). Реализованы
стандартные проверки корректности полей ввода: формат электронной
почты, минимальная длина пароля 8 символов, наличие имени
пользователя. Хеширование паролей выполнено алгоритмом bcrypt
с рабочим параметром $\text{cost} = 12$. После успешной
аутентификации клиент получает пару JWT-токенов в соответствии
со стандартом RFC~7519~\cite{rfc7519} и автоматически перенаправляет
пользователя в интерфейс, соответствующий его роли.

Для защиты от атак методом полного перебора перед эндпоинтами
аутентификации установлен ограничитель скорости запросов по IP-адресу
(token-bucket, не более 20 запросов в минуту), реализованный
в \texttt{internal/server/ratelimit.go}. Обновление токена доступа
выполняется по защищённому каналу при предъявлении валидного
refresh-токена; refresh-токены ротируются при каждом обновлении
в соответствии с рекомендациями OWASP ASVS~\cite{owasp_asvs}.
Принятые меры информационной безопасности перечислены
в табл.~\ref{tab:security-owasp} и соответствуют требованиям
ГОСТ Р 58833-2020~\cite{gost58833} и ГОСТ Р 56939-2024~\cite{gost56939}.

\begin{table}[!htb]
\fontsize{11pt}{13pt}\selectfont
\centering
\caption{Реализованные меры защиты по перечню OWASP Top 10}
\label{tab:security-owasp}
\smallskip
\begin{tabular}{|p{3.0cm}|p{10.5cm}|}
    \hline
    \textbf{Категория} & \textbf{Реализованная мера} \\
    \hline
        A01 Broken Access Control & Авторизация по принадлежности ресурса в каждом обработчике; дублирующие ограничения целостности БД \\
    \hline
        A02 Cryptographic Failures & bcrypt с $\text{cost}=12$; JWT с алгоритмом HS256; TLS 1.3 на канале клиент-сервер \\
    \hline
        A03 Injection            & Параметризованные SQL-запросы через драйвер \texttt{pgx} \\
    \hline
        A04 Insecure Design      & Принцип наименьших привилегий; идемпотентность операций записи через уникальные ограничения \\
    \hline
        A05 Security Misconfiguration & gin в release-mode; обязательная переменная окружения \texttt{JWT\_SECRET} \\
    \hline
        A07 Authentication Failures & Token-bucket ограничитель по IP; ротация refresh-токенов; политика паролей \\
    \hline
\end{tabular}
\end{table}

\subsection{Модуль преподавателя}

Главный экран преподавателя содержит сетку карточек принадлежащих
ему курсов (см.~Рисунок~\ref{fig:ui-teacher-courses} в Приложении~2).
Создание нового курса выполняется через форму, открываемую кнопкой
добавления. При выборе курса открывается экран его деталей со списком
связанных банков вопросов. Поддерживаются операции создания,
переименования и удаления банка, а также переход к редактору
вопросов и к экрану аналитики.

Редактор банка вопросов (см.~Рисунок~\ref{fig:ui-bank-editor})
обеспечивает создание, редактирование и удаление вопросов нескольких
типов: с одиночным выбором, с множественным выбором, с открытым
ответом. Реализован режим массового выделения для группового
удаления. Из редактора возможен запуск интерактивной сессии
и переход к аналитике. Реализована загрузка банков вопросов
из внешних источников~--- формат Moodle XML и Moodle GIFT,~---
что закрывает задачу импорта банка вопросов из системы дистанционного
обучения высшего учебного заведения, поставленную в разд.~1.3.3.

По завершении сессий преподавателю доступна психометрическая
аналитика по банку вопросов (см.~Рисунок~\ref{fig:ui-bank-analytics}),
реализующая модель оценки качества тестового инструмента
из подраздела~2.4.4: расчёт коэффициента $\alpha$ Кронбаха~\cite{cronbach1951},
индекса трудности $p_i$ и точечно-бисериальной корреляции $r_{pb,i}$
ответа на пункт с суммарным баллом для каждого пункта банка.
Результаты отображаются в форме психометрической карточки
с большим значением $\alpha$ и цветовой интерпретацией
(зелёная~--- $\alpha \geq 0{,}8$, синяя~--- $\alpha \geq 0{,}7$,
оранжевая~--- $\alpha \geq 0{,}6$, красная~--- $\alpha < 0{,}6$),
а также таблицы пунктов с цветовыми флажками: «слишком лёгкий»
при $p_i \geq 0{,}85$, «слишком сложный» при $p_i < 0{,}30$,
«низкая дискриминация» при $r_{pb,i} \leq 0{,}2$. Аналитика
позволяет преподавателю итеративно улучшать банк~--- удалять
неработающие пункты и переформулировать сомнительные. Расчёт
психометрических показателей инкапсулирован в чистой функции
\texttt{services.AnalyzeBank} (Приложение~1, листинг~П.1.4),
не зависящей от инфраструктуры и используемой как обработчиками
REST-запросов, так и в юнит-тестах.

\subsection{Модуль проведения интерактивной сессии}

Перед запуском сессии преподаватель задаёт её параметры: название,
режим проведения (\texttt{classic}, \texttt{timer} или \texttt{race}),
длительность таймера, признак перемешивания вопросов. По
подтверждению формируется новая запись в таблице сессий
и генерируется шестизначный числовой код для подключения участников.

Интерфейс преподавателя в ходе сессии отображает числовой код
подключения, текущий вопрос, агрегированную статистику ответов
и таблицу лидеров. Кнопка перехода к следующему вопросу переводит
сессию из состояния \texttt{active} в состояние \texttt{review},
а затем~--- к следующему вопросу. Интерфейс обучающегося
(см.~Рисунок~\ref{fig:ui-room-active} в Приложении~2) содержит
текущий вопрос и варианты ответа. После отправки ответа кнопки
блокируются; при переходе сессии в состояние \texttt{review}
подсвечивается правильный вариант. По завершении сессии участникам
отображается итоговая таблица лидеров
(см.~Рисунок~\ref{fig:ui-leaderboard}).

Архитектура канала WebSocket~\cite{rfc6455}, обеспечивающего
синхронизацию состояния сессии между клиентами, основана
на единственной горутине-реакторе, функционирующей без явных
блокировок. Все операции изменения состояния (регистрация клиента,
отключение, рассылка сообщений) проходят через каналы; клиенты,
не успевающие принимать сообщения, автоматически отключаются
от хаба для предотвращения блокировки рассылки остальным
участникам. Реализация хаба приведена в репозитории
\texttt{eduquiz-backend} (\texttt{internal/ws/}); типизированная
подписка на поток событий сессии на стороне клиентского приложения
реализована с использованием библиотеки
\texttt{web\_socket\_channel}.

\subsection{Модуль обучающегося и интервальное повторение}

Главный экран обучающегося (см.~Рисунок~\ref{fig:ui-student-home}
в Приложении~2) содержит блок отображения текущих очков опыта,
уровня и серии активных дней, а также поле ввода числового кода
интерактивной сессии. При наличии карточек, готовых к повторению
по алгоритму SM-2 (см.~подраздел~2.4.3), отображается уведомление
с указанием числа карточек к повторению на текущую дату.

Экран профиля обучающегося (см.~Рисунок~\ref{fig:ui-profile})
представляет агрегированную игровую статистику в рамках выбранного
курса: накопленные очки опыта, текущий уровень с индикатором
прогресса до следующего уровня, серию дней активности. График
активности за последние 14 дней отображает динамику накопления
очков опыта. Каталог достижений отображается сеткой плиток
с подсветкой полученных и описанием условий получения остальных.

Экран интервального повторения (см.~Рисунок~\ref{fig:ui-sm2})
предъявляет обучающемуся карточки, для которых наступила дата
следующего повторения, рассчитанная алгоритмом SM-2. После каждого
ответа сервер пересчитывает параметры состояния пары «обучающийся,
вопрос» по формулам~\eqref{eq:sm2-update} и~\eqref{eq:sm2-ef}
и планирует следующее предъявление. Оценка качества ответа $q$
формируется автоматически по правилу~\eqref{eq:q-score} из пары
«корректность~--- время», без явного запроса оценки у обучающегося.
Реализация алгоритма приведена в Приложении~1 (листинг~П.1.1);
алгоритм инкапсулирован в чистой функции без побочных эффектов,
что обеспечивает возможность модульного тестирования без подъёма
базы данных и переиспользование одной и той же реализации
как в интерактивной сессии, так и в режиме индивидуального
повторения.

\subsection{Инфраструктура развёртывания}

Контейнеризация сервиса выполнена с использованием
Docker~\cite{docker_docs} и Docker Compose. Одна composition
описывает как development-, так и production-среду с подменой
переменных окружения. В development-среде доступны два сценария
запуска: \texttt{make dev} запускает PostgreSQL и Redis
в Docker-контейнерах, а серверное приложение~--- локально
через \texttt{go run}; \texttt{make up} поднимает все сервисы
(включая серверное приложение) в контейнерах одной командой.

Production-развёртывание выполнено на виртуальной машине Yandex
Cloud с операционной системой Ubuntu Server LTS (2 vCPU, 2 ГБ RAM,
20 ГБ SSD). Обратный прокси Caddy~\cite{caddy_docs} обеспечивает
автоматический выпуск и обновление TLS-сертификата Let's Encrypt
по протоколу ACME. Деплой автоматизирован средствами GitHub Actions:
при push в ветку \texttt{deploy-demo} рабочий процесс выполняет
сборку Docker-образа на виртуальной машине, остановку и запуск
композиции и накат миграций схемы базы данных. Сервис доступен
по доменному имени \texttt{eduquizz.ru} с TLS 1.3.

\textbf{Расчётное обоснование показателей надёжности.} Архитектура
сервиса декомпозируется на пять узлов: серверное приложение,
PostgreSQL, Redis, Caddy и виртуальная машина хостинг-провайдера.
Для каждого узла приняты значения интенсивности отказов $\lambda_i$
и среднего времени восстановления $T_{\text{в}i}$ по эксплуатационным
данным аналогичных систем и соглашению об уровне обслуживания
провайдера в соответствии с подходом, изложенным в~\cite{ostreikovsky,gost27001}.
Для произвольного узла вероятность безотказной работы за время $t$
при экспоненциальном законе вычисляется по формуле $P_i(t) = e^{-\lambda_i t}$;
средняя наработка до отказа $\text{MTBF}_i = 1/\lambda_i$;
коэффициент готовности
$k_{\text{г}\,i} = \text{MTBF}_i / (\text{MTBF}_i + T_{\text{в}\,i})$.
В одиночной конфигурации без резервирования системный коэффициент
готовности рассчитывается как произведение по всем узлам
и составляет $k_\text{г}^\text{сист} \approx 0{,}9998$, что
соответствует ожидаемому времени недоступности около 1,75 часа в год.
Архитектурно предусмотрено горячее резервирование критичных узлов
(репликация PostgreSQL master-replica и горизонтальное масштабирование
stateless-приложения через балансировщик Caddy), что в расчёте
повышает системный коэффициент готовности до
$k_\text{г}^\text{сист*} \geq 0{,}9999$.

\subsection{Оценка производительности методом нагрузочного тестирования}

Для экспериментальной оценки производительности разработанного
сервиса реализован программный инструмент \texttt{cmd/loadtest}
на языке Go, выполняющий параллельную симуляцию заданного числа
обучающихся, одновременно проходящих сессию в режиме \texttt{race}.
Каждый виртуальный обучающийся последовательно выполняет операции
регистрации, подключения к сессии и ответов на все вопросы банка;
запросы виртуальных обучающихся параллелизуются через горутины.
По завершении прогона вычисляются перцентили p50, p95, p99 латентности
для каждого типа запроса, общее число и типы ошибок.

Тестовый стенд: рабочая станция Apple M-series; экземпляр сервера
EduQuiz, СУБД PostgreSQL~16 и сервер Redis~7 развёрнуты совместно
в Docker Compose. Прогоны выполнены при числе параллельных
обучающихся $N \in \{30, 100, 300, 500\}$ на банке из 10 вопросов.
Результаты приведены в табл.~\ref{tab:loadtest-results}.

\begin{table}[!htb]
\fontsize{11pt}{13pt}\selectfont
\centering
\caption{Перцентиль $p_{99}$ латентности эндпоинтов API EduQuiz
при $N$ одновременных обучающихся, мс}
\label{tab:loadtest-results}
\smallskip
\begin{tabular}{|l|c|c|c|c|}
    \hline
    \textbf{Эндпоинт} & $N=30$ & $N=100$ & $N=300$ & $N=500$ \\
    \hline
        \texttt{POST /auth/register} (bcrypt) & 1340 & 5975 & 15018 & 27191 \\
    \hline
        \texttt{POST /auth/login}             &  264 &  313 &   360 &   456 \\
    \hline
        \texttt{POST /rooms}                  &   11 &    6 &    42 &    28 \\
    \hline
        \texttt{POST /rooms/:id/start}        &    2 &    3 &    12 &    15 \\
    \hline
        \texttt{POST /rooms/join}             &   34 &   40 &    86 &   271 \\
    \hline
        \texttt{GET /rooms/:id/my/state}      &    7 &   25 &    67 &   120 \\
    \hline
        \texttt{POST /rooms/:id/my/answer}    &   12 &   34 &   142 &   233 \\
    \hline
        \textbf{Доля ошибок 5xx, \%}          & 0    & 0    & 0     & 0 \\
    \hline
        \textbf{Длительность сессии $N$ обучающихся, с} & 5{,}7 & 5{,}7 & 6{,}5 & 6{,}7 \\
    \hline
\end{tabular}
\end{table}

Анализ результатов показал, что для эндпоинтов критического пути
обработки запросов сессии (\texttt{POST /my/answer} и
\texttt{GET /my/state}) перцентиль $p_{99}$ латентности возрастает
сублинейно по отношению к числу обучающихся: с 12~мс при $N = 30$
до 233~мс при $N = 500$. В ходе проведённого нагрузочного тестирования
ошибки класса 5xx (внутренние ошибки сервера) не наблюдались.
Полученные значения позволяют заключить, что разработанная
архитектура в одиночной конфигурации без репликации устойчиво
обслуживает не менее 500 одновременных обучающихся в одной сессии,
что превышает реалистичный размер аудитории дисциплины
\textit{«Надёжность»} кафедры АСУ РГУ Нефти и Газа (НИУ)
имени И.М. Губкина в 5--10 раз. Тем самым выполнена задача оценки
производительности сервиса, поставленная в разд.~1.4.

Высокое значение латентности эндпоинта \texttt{POST /auth/register}
обусловлено вычислительной стоимостью алгоритма bcrypt с рабочим
параметром $\text{cost} = 12$ (порядка 250~мс на вычисление одного
хеша) и представляет собой обоснованный компромисс между скоростью
обработки запроса и стойкостью к атакам методом полного
перебора~\cite{owasp_asvs}. Эндпоинт регистрации не относится
к критическому пути обработки запросов сессии: регистрация
обучающегося выполняется однократно при первом обращении к сервису
и в реальной эксплуатации не совпадает по времени с моментом
начала учебной сессии. Соответственно, наблюдаемая на стенде
деградация перцентиля $p_{99}$ латентности регистрации при
$N = 500$ параллельных регистраций не препятствует обслуживанию
500 одновременных обучающихся в одной сессии и относится к режиму
синтетической нагрузки, не возникающему в учебном процессе.

\subsection{Методика педагогического эксперимента и пайплайн обработки}

Для последующей экспериментальной верификации педагогической
эффективности разработанного сервиса сформулирована методика
квазиэкспериментального исследования с контрольной (КГ)
и экспериментальной (ЭГ) группой. Гипотеза $H_1$: применение
сервиса EduQuiz с расширенной геймификацией и индивидуальным
интервальным повторением статистически значимо повышает успеваемость
и удержание знаний по сравнению с традиционным форматом обучения.

\textbf{Этапы исследования:} \emph{pre-test}~--- начальный замер
уровня знаний по теории надёжности до начала вмешательства;
\emph{обучающие сессии}~--- в ЭГ через сервис EduQuiz, в КГ ---
по традиционному формату занятий; \emph{post-test}~--- замер
уровня знаний по итогам цикла обучения; \emph{retention-test}~---
повторный замер через 1--2 недели для оценки долговременного
удержания знаний. Участие добровольное, данные обезличиваются
перед анализом, результаты исследования не влияют на текущую
успеваемость по курсу. Дисциплина и преподаватель: курс
\textit{«Надёжность»}, преподаватель~--- Волков Д.А., кафедра
автоматизированных систем управления; конкретные группы,
численность и даты тестов определяются по согласованию с научным
руководителем. Апробация на реальной группе обучающихся запланирована
в текущем семестре.

Для обработки результатов реализован программный пайплайн на языке
Python~3 с использованием библиотек \texttt{pandas},
\texttt{scipy}, \texttt{pingouin} и \texttt{matplotlib}. Этапы
обработки:
\begin{itemize}
    \item проверка нормальности распределений критерием
          Шапиро\textendash Уилка ($p \geq 0{,}05$);
    \item сравнение средних между КГ и ЭГ: критерий Манна\textendash Уитни
          при отклонении от нормальности, критерий Стьюдента/Уэлча
          при выполнении предпосылок нормальности;
    \item оценка размера эффекта по коэффициенту Коэна $d$
          с непараметрической оценкой доверительных интервалов;
    \item поправка Бонферрони при множественных сравнениях
          трёх замеров ($\alpha' = 0{,}05 / 3$);
    \item расчёт коэффициента $\alpha$ Кронбаха~\eqref{eq:cronbach}
          и точечно-бисериальной корреляции $r_{pb}$
          для каждого пункта тестового инструмента.
\end{itemize}

Корректность пайплайна проверена на модельных данных, отражающих
типичные паттерны педагогических экспериментов. Применение пайплайна
к реальным данным выполняется после завершения цикла обучающих
сессий путём замены входных файлов без изменения кода пайплайна.

\subsection{Перспективы дальнейшей научно-исследовательской работы}
\label{sec:future-work}

К моменту завершения работы сервис EduQuiz реализован как прототип,
обеспечивающий основные сценарии использования преподавателя
и обучающегося; на его основе формируется задел для дальнейшей
научно-исследовательской работы. Перспективы расширения сервиса
и направления научно-исследовательской работы включают следующее.

\textbf{Полная интеграция с системой дистанционного обучения
университета.} В реализованной версии загрузка банков вопросов
выполняется в форматах Moodle XML и GIFT. Полная двусторонняя
интеграция с порталом edu.gubkin (синхронизация курсов и списков
обучающихся, отправка результатов сессий в журнал) запланирована
в дорожной карте развития.

\textbf{Промышленный мониторинг и резервирование.} Подключение
системы мониторинга Prometheus с экспортом per-endpoint латентности
и доли ошибок, визуализация в Grafana и алертинг через Alertmanager.
Репликация PostgreSQL master-replica с автоматическим failover
для повышения коэффициента готовности; целевое RTO 30 секунд,
RPO 1 секунда.

\textbf{Адаптивное тестирование на основе теории отклика на пункт.}
Ключевая исследовательская проблема, выделенная по результатам
работы, состоит в том, что существующие интерактивные тестирующие
системы измеряют формальную правильность ответа на отдельный
вопрос, но не оценивают латентный уровень подготовки обучающегося
независимо от состава конкретного теста. Совместное моделирование
параметров банка вопросов методами теории отклика на пункт
(Item Response Theory, IRT~\cite{crocker_algina}) и оценка
латентного уровня $\hat{\theta}$ обучающегося, а также построение
адаптивного режима тестирования с подбором следующего вопроса
по критерию максимума функции информации Фишера в окрестности
текущей оценки $\hat{\theta}$, образуют содержание дальнейшей
научно-исследовательской работы. Реализованные в настоящей работе
архитектура сервиса, модель данных и психометрическая аналитика
средствами классической теории тестов формируют достаточную основу
для реализации этих направлений.

\textbf{Оценка валидности атрибуции ответа в условиях распространения
больших языковых моделей.} Дополнительной исследовательской задачей,
получающей актуальность в условиях широкого распространения
общедоступных сервисов на основе больших языковых моделей
(ChatGPT, GigaChat, YandexGPT и др.), является разработка модели
оценки достоверности самостоятельного выполнения учебного
тестирования обучающимся. На уровне модели данных в разрабатываемом
сервисе предусмотрено хранение поведенческой телеметрии каждой
попытки ответа (время ответа, последовательность переключений
активного окна, идентификатор устройства), необходимой для
последующего построения классификатора валидности атрибуции
попытки. Разработка такого классификатора и интегральной оценки
знаний, учитывающей одновременно измерительную точность тестового
инструмента и валидность атрибуции ответа, выделена
как самостоятельное направление последующей научно-исследовательской
работы.

\textbf{Выводы по главе 3.}
В главе описаны четыре реализованных и работающих в режиме
«end-to-end» функциональных модуля сервиса EduQuiz: модуль
аутентификации с реализацией мер защиты по перечню OWASP Top 10
и требованиям ГОСТ Р 58833-2020 и ГОСТ Р 56939-2024; модуль
преподавателя с управлением курсами и банками вопросов и реализованной
психометрической аналитикой (коэффициент $\alpha$ Кронбаха, индекс
трудности, точечно-бисериальная корреляция); модуль проведения
интерактивной учебной сессии в трёх режимах с синхронизацией через
WebSocket-канал; модуль обучающегося с индивидуальным интервальным
повторением по алгоритму SuperMemo SM-2 и игровой прогрессией.
Сервис развёрнут на доменном имени \texttt{eduquizz.ru}
с автоматическим выпуском TLS-сертификата. Методом нагрузочного
тестирования установлено, что разработанная архитектура устойчиво
обслуживает не менее 500 одновременных обучающихся в одной сессии
при $p_{99}$ латентности критического пути не выше 233~мс и нулевой
доле ошибок класса 5xx; нефункциональное требование к производительности
выполнено. Разработана методика педагогического эксперимента
и реализован программный пайплайн её обработки методами
математической статистики; корректность пайплайна проверена
на модельных данных, апробация на реальной группе обучающихся
запланирована в текущем семестре. Сформулированы перспективы
дальнейшей научно-исследовательской работы~--- задел для
последующих этапов исследования.

\newpage
